\scanpage{120}
% archive sheet number: 119
Plaçons nous pour fixer les idées sur un corps de base, et voyons \unsure{ds.|dans|sous} quelles conditions de telles données définissent bien un polyèdre régulier \inserted{\unsure{inscrit}}. Tout d'abord il faut que le noyau
\[
 N=\mathrm{Ker}\,\unsure{\varphi|\rho}\ \subset\Gamma\tag*{(30)}
\]
soit d'\uline{indice fini}. Ceci étant, la condition \deleted{\unsure{la c.}} (17) est triviale par construction, mais la condition $\boxed{11}$ semble plus sérieuse, et s'écrit ceci, en posant
\[
 s'_0=\sigma_0(s_0),\quad a'_0=\sigma_1(a_0),\quad f'_0=\sigma_2(f_0)\tag*{(31)}
\]
$\boxed{32}$ \uline{On a} $s'_0\ne s_0,\ a'_0\ne a_0,\ f'_0\ne f_0$, et
\[
 \begin{aligned}
 \text{a)}\quad&\text{Si }\gamma\in\Gamma\text{ est tel que }\gamma s_0\unsure{\prec|<}a_0,\ \text{alors }\gamma s_0\in\{s_0,s'_0\}\\
 \text{b)}\quad&\text{Si }\gamma\in\Gamma\text{ est tel que }s_0\prec\gamma a_0\prec f_0,\ \text{alors }\unsure{\sigma|\gamma}a_0\in\{a_0,a'_0\}\\
 \text{c)}\quad&\text{Si }\gamma\in\Gamma\quad\text{---}\quad\deleted{\unsure{\gamma}}a_0\prec\gamma f_0,\ \text{alors }\gamma f_0\in\{f_0,f'_0\}
 \end{aligned}
\]
\[
 \begin{aligned}
 \gamma s_0\prec\gamma'a_0\prec\gamma''f_0&\ \Longrightarrow\ \exists\gamma'''\in\Gamma,\ \text{tel que}\\
 \gamma'''s_0=\gamma s_0,\quad\gamma'''a_0&=\gamma'a_0,\quad\deleted{\unsure{\gamma'''}}\gamma'''f_0=\gamma''f_0
 \end{aligned}\tag*{\boxed{33}}
\]
(exprime $\boxed{\unsure{19|17}}$)

% relocation arrow: block written above (33) is directed to this position below (33).
\inserted{De plus $\boxed{18}$ s'écrit : $\boxed{\unsure{34|35}}$ les \deleted{\unsure{les}} $\gamma s_0$ engendrent affinement $X$ (première condition !)}
